\documentclass[10pt,a4paper]{amsart}
\usepackage[margin=19mm]{geometry}
\usepackage{amsmath,amssymb,mathtools}
\usepackage[scr=rsfs,cal=euler]{mathalfa}
\usepackage{xfrac}
\usepackage[unicode,hidelinks,bookmarks=false]{hyperref}
\newcommand{\ie}{i.e.\ }
\newcommand{\cf}{cf.\ }
\newcommand{\resp}{respectively\ }
\newcommand{\Subsection}{\subsection}
\providecommand{\nobreakdash}{\nobreak\hbox{-}}
\setlength{\emergencystretch}{2em}
\multlinegap0pt
\usepackage{amssymb}
\usepackage{mathtools}
\usepackage{bm}
\usepackage{xcolor}
\usepackage[shortlabels]{enumitem}
\usepackage{tikz}
\newtheorem{proposition}{Proposition}
\newcommand{\cV}{\mathcal{V}}
\renewcommand{\tilde}{\widetilde}
\title{Real Geometric transcendence for the Gamma function}
\author{Arshay Sheth}
\date{Source: arXiv:2605.15117v1 (CC BY 4.0)}
\begin{document}
\maketitle

\noindent\emph{Source and licence.} Real Geometric transcendence for the Gamma function. Version arXiv:2605.15117v1 (CC BY 4.0), distributed by arXiv under the Creative Commons Attribution 4.0 International licence, http://creativecommons.org/licenses/by/4.0/, as recorded in the arXiv licence metadata for this exact version and shown on the abstract page https://arxiv.org/abs/2605.15117v1. Full licence notice: \texttt{licenses/arxiv-2605.15117-CC-BY-4.0.txt}.\par\smallskip

\noindent\textit{Editorial scope.} Counted unit: the article's proposition lem:realcompbalg (source0.tex lines 371--373). The article's complete original proof of it is delivered with it (source0.tex lines 374--379, one \texttt{proof} environment of 6 lines). No mathematics is written, repaired or re-derived: the statement and the proof are the article's own lines, in the article's own notation, and no retained line is substituted or re-flowed.  No further source-local context is needed: every cross-reference of every retained line resolves inside the two delivered spans.  Nothing was omitted from any retained span, so the package carries no omission record.  No retained line contains a citation, so the wrapper carries no bibliography and no bibliography entry.  No retained line contains a figure environment, an image-inclusion command or drawing code, so no image is delivered and the package declares no asset dependency.  Editorial formatting only, in addition: an AMS \texttt{amsart} preamble that restates the article's own declarations and macros used by the retained lines, namely the environments \texttt{proposition} on separate counters, together with the standard packages those lines need.  The retained environments print in this document as Proposition 1 in order of appearance, whereas the article prints the same items with the numbering of its own full document.  The complete native source of the article as distributed in the arXiv e-print for this exact version is bundled unchanged as \texttt{sources/200634/source0.tex}.
\begin{proposition} \label{lem:realcompbalg}
The curve $f(C_P)$ is bialgebraic for $\Gamma^2$.
\end{proposition}

\begin{proof}
 Since $\mathcal V$ is infinite, $\tilde{\cV}$ contains a subset $I$ which is homeomorphic to an open interval.  Since any complex curve has only finitely many singular points, we may choose a smooth point $x$ of $C_P$ lying in $I$.  By the implicit function theorem, there exists an open subset $U \subseteq C_P$ (in the complex analytic topology) containing $x$, an open disc $D \subseteq \mathbb C$,  and a biholomorphic map $\phi: D \rightarrow U$.  By shrinking $U$ if necessary we  may also assume that  $f(U) \subseteq (\mathbb C  \smallsetminus \mathbb Z_{\leq 0})^{2}$. We now note that $(Q \circ g \circ q) \circ \phi$ vanishes on $\phi^{-1}(I \cap U)$; by the identity theorem, this function must vanish on $D$ as well.  So $Q \circ g \circ q$ vanishes on $U$ and hence $U \subseteq A$. 
This implies that the analytic subset $A$ of $C_P$ is not thin,  
%Since $C_P$ is irreducible, a standard result implies that it is  connected in the Euclidean topology 
and so we must have $A=C_P$. Therefore $q(C_P)\subseteq \hat{C}_Q$ and so $\Gamma^2(f(C_P)) \subseteq \hat{C}_Q $ which in turn implies that $\Gamma^2(f(C_P) \cap (\mathbb C  \smallsetminus \mathbb Z_{\leq 0})^{2} ) \subseteq g^{-1}(C_Q)$, proving that  $f(C_P)$ is bialgebraic for $\Gamma^2$. 
\end{proof}

\end{document}
